

\section{Threat Model}

\revised{
Given people's unawareness of the security risk of mini-app interaction history, the next part of this paper will reveal the dangers. In this section, we will first introduce the threat model in which the interaction data can be used to threaten users' privacy.
}

\noindent \textbf{Scenario.}
We consider a leading global super-app vendor that consists of multiple departments, each with distinct roles and responsibilities. 
In this scenario, there are two types of departments: the high-privilege business department ($D_{high\_priv}$), such as payment and risk control, and the low-privilege business department ($D_{low\_priv}$), such as lifestyle or third party-collaborators. Both  $D_{high\_priv}$ and $D_{low\_priv}$ strictly comply with privacy protection measures to ensure service security and regulatory compliance. $D_{high\_priv}$ has access to sensitive data, such as personal financial information, while $D_{low\_priv}$ can only access data explicitly classified as non-sensitive by the super-app vendor and should know minimal user privacy. 

\revised{In this scenario, the insider threats~\cite{6957307} within the $D_{low\_priv}$ may act as the adversary. These insiders cannot access private data but can access mini-app interaction history because it is considered non-sensitive. For instance, employees in advertising or loan departments cannot directly access sensitive user attributes due to data-isolation policies. However, they can access \mini and \op to infer those attributes indirectly. Their motivation may include selling the inferred data for profit and gaining a competitive advantage by enhancing targeted advertising or differential pricing strategy. Even when their mini-app does not hold explicit data (e.g., a user’s assets status), relevant privacy attributes of unlabeled users can still be inferred by analyzing broader interaction patterns. As an example, identifying whether a user owns property or a vehicle helps loan department insiders tailor services more effectively. Meanwhile, since users rarely access the loan department’s mini-app, the department cannot gather enough privacy attributes directly. Nonetheless, these attributes can still be inferred by analyzing the mini-app interaction history.}


\noindent \textbf{Adversary's Goal.}
\revised{The adversary's goal is to \textit{accurately predict the privacy attributes of as many users as possible from the non-sensitive data}.}
The goal consists of twofold. 
% 
First, the adversary aims to identify a user subset whose privacy attributes are strongly correlated with \mini and \op. These users are believed to have sufficient mini-app interaction history data to
infer their privacy attributes. 
For example, one user might frequently use the fueling mini-apps, which could indicate that he/she owns a vehicle.
For other users, if the \mini and \op cannot reveal their privacy attributes, the adversary may classify them as unknown rather than forcibly assign labels. This can avoid blind and inaccurate predictions.
% 
Second, for users in the subset, the adversary aims to infer their privacy attributes as accurately as possible. This is because privacy attributes are valuable but scarce. With such attributes, adversaries can design unfair market strategies (e.g., targeted advertising, price discrimination, or denial of services) for certain user groups. For example, they might use financial status to implement differential pricing that targets wealthier users with higher service fees. Meanwhile, all super-app users are potential victims because many users do not provide such information due to privacy concerns. For instance, according to \alipay's data, only 6.2\%, 5.8\%, and 6.9\% users provide their marital status, property ownership, and vehicle ownership, respectively. 




Note that this goal is
realistic and dangerous because super-app vendors often have a large
user base (over one billion). Even if the adversary can only infer the privacy of
a small portion of users, the number of threatened users is significant. Although the super-app vendor strives to control access by $D_{low\_priv}$, as super-app functionality expands, data initially considered privacy-irrelevant may become more privacy-revealing, leading to new risks.


\noindent \textbf{Adversary's Ability.}
\revised{Unlike prior work~\cite{tu2018your} that requests a vast number of OS-level data, we assume the adversary can acquire only \mini and
\op from the vendor's server because existing super-apps (\eg Alipay, WeChat, and Tiktok) are web-based and hosted on the servers~\cite{alipaysdk,wechatsdk,tiktoksdk}, the vendor naturally needs \mini and \op to provide the desired app functionality.} In practice, the privilege level of different data is determined by regulations, industrial practices, or research papers. Since \mini and \op are not identified as sensitive in these above sources, as discussed in Section~\ref{sec:privacy_of_history}, they are not identified as sensitive and can be accessed by $D_{low\_priv}$.



Due to regulation and performance constraints, we assume the attacker cannot modify the released versions of the super-apps on users' devices, such as instrumentation or probing, to collect user information on the device side.


We also assume that the adversary can collect the privacy attributes of a small number of users from a publicly available data breach. These breaches are widely available on the Internet~\cite{liu2015cloudy,mayer2021now,thomas2017data}, which contains phone numbers, personal identification numbers, and other sensitive privacy attributes. The adversary can match their user data to the samples in breach databases (e.g., by matching phone numbers~\cite{cheng2017enterprise, thomas2017data}) to identify the privacy attributes of a few hundred users. The adversary can then collect \mini and \op of these users to train the model to infer the privacy of other users.

\section{THEFT: Mini-app Interaction History-Oriented Inference Attack}

\revised{
To reveal the new security risk of mini-app, we design a new attack, \tool, \textit{mini-app interac\underline{T}ion \underline{H}istory-ori\underline{E}nted in\underline{F}erence a\underline{T}tack}. This attack is proof of the potential security issue of the interaction data.}
\tool uses both \mini and \op to infer the privacy attributes of users. We will first illustrate the overall pipeline, followed by a detailed presentation of each component. 





\subsection{Overview}


Generally, the \tool employs a DNN model to infer user privacy. The inputs of the DNN model are \mini and \op, and the model outputs are the predicted attribute label and a confidence score, which is used to filter out the high-confidence victim subset. The confidence score needs to approximate the real accuracy of the prediction. For example, if a model predicts the user to be a female with a confidence level above 0.9, then the accuracy of this prediction should be more than 90\%. Using a predefined threshold, the adversary can select high-confident samples to identify a vulnerable victim subset and provide label predictions while labeling other users with lower confidence as unknown.


The attack process consists of three stages, as shown in
Figure~\ref{fig:attack}. The first step is the attack model training, where the
adversary uses the \mini and \op data of a relatively small set of leaked users with one-hot privacy attribute labels to train a DNN
model.
% 
The second step is model confidence calibration, where the adversary uses a
calibration technique to ensure that the confidence score can faithfully reflect
the accuracy of the prediction.
% 
The last step is online inference, where the adversary uses the trained
model to predict the privacy attributes of unlabeled users from \mini and \op.




\begin{figure*}[!t]
    \centering
    \includegraphics[width=0.95\linewidth]{fig-empirical-attack.pdf}
    \vspace{-1em}
	\caption{The pipeline of \tool consists of three steps: attack model training, confidence calibration, and online inference.}
	\label{fig:attack}
 \vspace{-1em}
\end{figure*}


\subsection{Attack Model Training}


\begin{table}[!tb]
\centering
\caption{28 categories of mini-apps.}
\vspace{-1em}
\label{tab:category}
\begin{adjustbox}{max width = 1.\columnwidth}
\begin{tabular}{@{}cl|cl@{}}
\hline
\textbf{Code} & \textbf{Category} & \textbf{Code} & \textbf{Category} \\ \hline
1 & \texttt{Education} & 15 & \texttt{Finance} \\
2 & \texttt{Entertainment} & 16 & \texttt{Food and Drink} \\
3 & \texttt{House and Home} & 17 & \texttt{Health and Fitness} \\
4 & \texttt{Lifestyle} & 18 & \texttt{Art and Design} \\
5 & \texttt{Maps, Navigation, and Taxi} & 19 & \texttt{Books} \\
6 & \texttt{Music and Audio} & 20 & \texttt{Comics} \\
7 & \texttt{Parenting} & 21 & \texttt{Communication} \\
8 & \texttt{Shopping} & 22 & \texttt{Medical} \\
9 & \texttt{Auto and Vehicles} & 23 & \texttt{News} \\
10 & \texttt{Beauty} & 24 & \texttt{Photo} \\
11 & \texttt{Business} & 25 & \texttt{Productivity} \\
12 & \texttt{Dating} & 26 & \texttt{Sports} \\
13 & \texttt{Social} & 27 & \texttt{Weather} \\
14 & \texttt{Travel and Local} & 28 & \texttt{Event} \\ \hline
\end{tabular}
\end{adjustbox}
\vspace{-1em}
\end{table}

In the attack model training stage, \tool trains a model that can better
learn the features of mini-app interaction history data and thus can accurately
predict the privacy attributes of users. The training data are the mini-app
interaction history (\mini and \op), the training labels are
collected privacy attributes from data breaches. We select the
transformer-based~\cite{vaswani2017attention} model as the default architecture
of \tool. The transformer model comprises 12 encoders, each equipped with
12 bidirectional self-attention heads, resulting in a total of 110 million parameters. In Section~\ref{sec:architecture_comparison}, we also compare our transformer
architecture with another two architectures to demonstrate the superiority of
our choice. For each type of privacy attribute, we modified the output dimension of
the last fully connected layer to the number of labels.


\noindent \textbf{Model Input.}
The model input consists of two parts, \mini (denoted as $x_m$) and \op
(denoted as $x_o$). $x_m$ contains the list of the latest $N$ mini-apps accessed
by a user. For the $i$-th mini-app in the list, we record 1) the unique id
$m^{i}_{id}$ of the mini-app (maintained by the \alipay backend), 2) the
mini-app category code $m^i_{c}$, and 3) the number of access times $m^i_{f}$
over the last 30 days. For the category code of mini-apps, we classify the
mini-apps into 28 types following prior literature~\cite{li2022collect}. Table~\ref{tab:category} displays all the categories. Therefore, $x_m$ is a tensor
of the shape of $(M, 3)$. 
Meanwhile, $x_o$ contains button IDs that the user clicked over the past $M$ timestamps, where each timestamp represents a 500-millisecond interval. \revised{We convert the original click logs with timestamps from \alipay into this format, to unify the input format for our deep learning model.}
Within each interval, if the
user clicks a button in the mini-app, we record its ID. Otherwise,
we record $0$ to indicate that the user does not click any button. These button IDs are globally unique in the super-app, with each ID corresponding to a specific function and remaining consistent across all mini-apps. The input of the model is the fusion of
$x_m$ and $x_o$. To fuse the two types of data, we empirically set $N=M$ and
concatenate $x_m$ and $x_o$ into a single input tensor for the model.



\noindent \textbf{Model Output.}
The output is the predicted label for a specific type of privacy attribute. In this
paper, we focus on seven privacy attributes: gender, location, age, property
ownership, vehicle ownership, marital status, and parental status. All these
types of privacy are maintained by \alipay and are considered important privacy
information in the internal company. For each privacy attribute, we regard it
as a classification task. For five privacy attributes, the inference task is a
binary classification task. The five privacy attributes are gender, whose labels are
male and female, and the ownership of property and vehicles, marital and
parental status, whose labels are yes and no. For location, we give three
labels: Tier-1 cities, Tier-2 cities, and Tier-3 cities, making this task a ternary
classification task. \revised{Specifically, Tier-1 cities include municipalities, well-developed provincial capitals, and economic centers; Tier-2 cities encompass other provincial capitals and major cities; and Tier-3 cities comprise ordinary cities and regions primarily composed of towns and villages.} Regarding age, we follow prior research and assign four labels:
Under 18, 18$\sim$39, 40$\sim$65, and Above 65~\cite{shauly2020public}.




\subsection{Model Confidence Calibration} 

\textbf{Necessity.}
Recall that the goal of \tool is to identify a subset of users whose
privacy attributes can be inferred accurately. However, the original
output confidence of the model cannot faithfully reflect the
inference accuracy. This is because the original model is optimized by the
cross-entropy (CE) loss. The CE loss will lead to overconfidence because it greedily maximizes
the confidence of the predicted label. However, for samples in which the model
is not confident, we do not want to give a high confidence score to
the predicted label because it will lead to a high false positive rate. In other
words, when the attack model generates a confidence score of 0.9 for the
prediction of user privacy, it should also ensure that the prediction accuracy
is above 0.9. Therefore, by selecting predictions with high confidence scores,
we can identify the subset of users from which we can accurately infer their
privacy attributes.


\noindent \textbf{Temperature Calibration.}
We adopt model calibration techniques~\cite{muller2019does,wang2021rethinking}
to achieve our goal. The general idea of model calibration is to append
an extra trainable calibration module after DNN's output layer. 
The calibration module takes DNN's output (overconfident
prediction) as input and generates a new confidence score. The calibration
module is fine-tuned on a calibration dataset to minimize the gap between the
generated confidence score and the real accuracy. Our design chooses the
temperature calibration~\cite{muller2019does,wang2023tabi} because it performs best in our evaluation. In Section~\ref{sec:ablation_study}, we
compare our temperature calibration with two other calibration techniques to
demonstrate the superiority of our choice.



Specifically, given a trained attack model, we replace its softmax function with
a calibrated softmax function, which is shown in Equation~\ref{eq:cali_tem}. The
calibrated softmax adds a trainable parameter, temperature $t$, to the softmax
function. Then we fine-tune $t$ by minimizing the negative log-likelihood loss
(NLL loss) in the validation dataset. Note that we fixed the parameters of the
attack model during the calibration phase and only updated $t$. The loss function choice and the validation dataset's use are consistent
with prior work~\cite{pmlr-v80-kuleshov18a,wang2023tabi}. 
\begin{equation}
    C(x) =\frac{\max_{i \in K}\exp(o_i / t)}{\sum_{i \in K} \exp(o_i / t)}
    \label{eq:cali_tem}
\end{equation}




\subsection{Online Inference} 

For inference, we use a pre-defined threshold $t_d$ as the confidence bar to
determine whether our model can confidently infer the privacy attribute of a
user. For the input of a user $x$, the output confidence score of the calibrated
module is $C(x)$. \tool returns ``unknown'' if $C(x)$ is smaller than $t_d$.
Otherwise, \tool will generate the corresponding attribute label for the given user.



\section{Evaluation}


In this section, we comprehensively evaluate the inference performance of \tool. 
We first describe the experimental setup in Section~\ref{sec:exp_setup}. 
In Section~\ref{sec:attack_effectiveness}, we
display the effectiveness of \tool, particularly the high inference accuracy
on the high-confidence samples. After that, in Section~\ref{sec:exp_insights},
we provide insight and in-depth analysis on the success
of \tool for each of the privacy attributes. In Section~\ref{sec:architecture_comparison}, we prove the effectiveness of our model architecture selection.
In Section~\ref{sec:ablation_study}, we conduct ablation studies to
demonstrate the generalizability of our approach and the recommended parameter
settings. 


\noindent\revised{\textbf{Ethical Disclaimer. }Since our evaluation involves collecting privacy attributes and mini-app interaction history, we obtained approval from the IRB of \alipay. For details, please refer to the Ethics Considerations Section.} 


\subsection{Experiment Setup}
\label{sec:exp_setup}


\textbf{Volunteer Selection.}
We first select 5\% of daily active users of \alipay as volunteer candidates and
send invitations to them. To guarantee the quality of the data, we asked the
internal engineers of \alipay to confirm that the candidates have been active in
\alipay for more than 7 days at the time of invitation. In each invitation, we
provided a clear consent notice on what information is collected and stated
that the data were collected only for research purposes. The candidates who
accepted the invitation become our participants. In total, we selected
38,351,206 candidates, of which 288,895 users agreed to share their data. 


For our participants, after their consent, we pushed the \alipay-dev version to their devices, replacing the original version. The
\alipay-dev version is the same as the original version. The only difference is that it displays an additional user agreement when the user first opens the app. This agreement notifies the users that he/she is involved in our experiment and displays detailed information about our experiments. Note that we do not add extra instrumentation or probes in \alipay-dev. We store \mini and \op in an encrypted database to protect the privacy of volunteers.



\noindent \textbf{Data Collection.}
We collect data during the whole process of using \alipay. Specifically, the
collection phase starts when the user opens the app and ends when the user exits
\alipay or switches to another app. We collect both \mini and \op during this phase. We collected data for 50 days and obtained 10,987,662 data samples from 288,895 users. 



We empirically set the length $N$ to 200 for each mini-app interaction history data sample. Specifically, for \mini, we record the Top 200 recently used mini-apps. We confirmed with the internal engineers of \alipay that 200 mini-apps
are sufficient to represent the user's recent behavioral patterns, as the
average number of mini-apps used by \alipay users is 117.8.
\revised{For \op, we convert the original logs in \alipay with timestamps into 0.5-second intervals over a 100-second window, resulting in a 200-dimensional vector per data sample.}
We empirically determined these values based on \alipay's report that the average time of user interactions is 89.2 seconds. A period of 100 seconds is sufficient to cover most
user interactions. The internal engineers of \alipay also confirmed that the
500-millisecond interval is the minimum for user interactions, and most
interactions are shorter than 100 seconds. Recent literature also supports the setting of these values~\cite{10.1145/2750858.2805837,10.1145/3604271}. Therefore, the
dimension of each mini-app interaction history data sample is $200 \times 4$. 



\noindent \textbf{Building the Ground Truth.} 
We strictly adhered to IRB requirements (illustrated in the Ethics Considerations
Section) to build ground truth privacy attributes as training labels.
For the location attribute, we obtained user consent to access the data of the system
location service. For other attributes, we first ask users to provide their
information. Then we sent authorization requests to acquire user consent to
access users' data in \alipay's server, thereby obtaining
gender, age, as well as property and vehicle information registered under their
names to confirm user-provided labels. 
For users who are unwilling to share all
their privacy attributes or the shared attributes do not match the data in \alipay's server,
we removed their data in our evaluation. In summary, we construct the dataset with 1,099,130 data samples from 219,826 users.


\noindent\textbf{\revised{Unbiasedness of Collected Dataset.}}
\revised{We further demonstrate that the collected dataset closely aligns with real-world population distributions.}
We analyzed label distribution for various privacy attributes: 22.8\% users are from tier-1 cities, 43.4\% from tier-2, and 33.8\% from tier-3. The gender distribution is almost equal: 50.3\% male and 49.7\% female. For age, 17.5\% are under 18, 38.0\% are between 18-40, 30.9\% are 40-65, and 13.6\% are above 65. Moreover, 48.6\% users own property, 21.0\% users own vehicles, 69.7\% are married, and 53.6\% have children. To verify the unbiasedness of our data, we conducted a chi-square test comparing these distributions with data from China’s National Bureau of Statistics~\cite{NBSC}. \revised{The test confirmed no significant statistical difference, indicating that our dataset is indeed representative of the wider population.}

\noindent \textbf{Experiment Protocol.} 
We split the collected data into three sets: a training set, a validation set, and a test set. To demonstrate that \tool can use a very small dataset as the training set to infer the privacy attributes of large-scale users, we set the size of training and validation sets to a small number (less than 0.1\% of the test set). Specifically, the training set includes 200 users, and the validation set includes another 200 users. We collect five data samples from each user. Thus, the total number of training and validation sets is 1,000. The test set contains 1,097,130 data samples from the other 219,426 users.
Furthermore, we repeated the evaluation ten times with random data partitioning to provide a reliable assessment.
For all comparisons in our evaluation, we run the hypothesis test~\cite{chen2023obsan} to ensure the significance of our observation.




\noindent \textbf{Implementation Details}
We conducted our experiment on a server with an Intel Xeon E5-2678 v3 CPU (48
cores), 128GB RAM, and 2 NVIDIA RTX 3090 GPUs. The server OS is Ubuntu 20.04
LTS. We implemented our code base using Python 3.8 and PyTorch 2.0. All models are trained for 100 epochs. 
We monitor the loss in the validation set and ensure that the models are converged.
We set the discriminator threshold ($t_d$) to 0.9, meaning we identify data samples with a confidence score greater than 90.0\% as a high-confidence subset, according to the literature~\cite{kuppers2020multivariate,wang2021confident}. 
We utilize Adam as the optimizer, and the learning rate is set to 1e-3. 
%\vspace{-1em}

\subsection{Attack Effectiveness}
\label{sec:attack_effectiveness}

In this section, we will first introduce the general effectiveness of \tool. Then, we will show that
\tool can effectively select a subset of users and accurately infer their
privacy attributes. Finally, we will illustrate that output confidence
of our model is aligned with the true accuracy of the prediction,
which is an essential indicator of the attack's effectiveness.

\begin{figure*}[!t]
    \centering
    \includegraphics[width=1.\linewidth]{fig-overall.pdf}
    \vspace{-1.5em}
	\caption{Attack effectiveness of \tool. (a) Overall inference accuracy for each type of privacy. (b)-(h) For each privacy type and confidence interval, the proportion of users predicted to be within the interval ($P_{int}$), the proportion of correctly predicted users ($P_{conf}$), and inference accuracy ($acc_{int} = P_{conf} / P_{int} \times 100\%$).}
	\label{fig:results_all}
 \vspace{-1em}
\end{figure*}


\noindent \textbf{General Effectiveness.}
\revised{In Figure~\ref{fig:results_all}~(a) we show the overall inference accuracy across all inference samples for
each privacy attribute label.} 
We ran each experiment 10 times and reported the average
 accuracy. 
We manually checked the standard error and confirmed that for all
labels, the standard error is below $0.7\%$. Therefore, we omit the standard
error in the figure. In Figure~\ref{fig:results_all}~(a), we also mark the accuracy of the random
guess baseline with a dashed line. The average inference accuracy of \tool
is $65.2\%$, which is $48.1\%$ higher than the baseline. In particular, for the age of users, \tool is $2.9\times$ higher than the baseline. 
We also performed a chi-square test to confirm the superiority of \tool. Our
null hypothesis is that the performance of the model is equal to random guessing, and the $p$ value is $0.02$. Therefore, we reject the null hypothesis ($p < 0.05$) and conclude that \tool is significantly better
than the baseline.



\noindent \textbf{Performance on High-Confident Subset.}
In Figure~\ref{fig:results_all} (b)-(h), we split the model confidence score
into intervals of 10\%. Each subfigure represents the results of one type of
privacy. For each confidence interval, we show the proportion of samples within
this interval ($P_{int}$, represented in light blue bars) and the
correctly predicted samples within that interval ($P_{conf}$, represented in
dark blue bars). Note that the more dark blue bars cover the light blue bars,
the more samples of $P_{int}$ are correctly predicted (a higher
attack performance). We also show the prediction accuracy within that interval
($acc_{int} = P_{conf} / P_{int} \times 100\%$, represented by the red line).
For each subfigure, the y-axis on the left represents the prediction accuracy
($acc_{int}$), and the right y-axis represents the proportion of samples within
the interval ($P_{int}$ and $P_{conf}$). 


For samples predicted by our model with high confidence scores, the
attack performance is \textit{near-perfect}. Setting the confidence
threshold as 90\%, we can on average identify 16.1\% samples, and the accuracy of the inference is 95.5\%. For the bar at the index 90-100 in
Figure~\ref{fig:results_all}~(b)-(h), the dark blue bars cover almost entirely the light blue bars. which means that the proportion of correctly predicted samples
($P_{conf}$) is very close to all samples in this interval ($P_{int}$). This
implies that our attacks are very likely to be successful in 15.4\% of all data
samples. 


We can also observe that high attack performance is consistent across all
types of privacy. For six out of seven types of privacy, we can identify more than 10\% users with an accuracy of more than 90\%. For gender, we can identify
11.0\% data samples with an accuracy of 97.7\% by setting the confidence
threshold to 90\%. For location and age, we can identify 12.0\% and 23.9\% data
samples with an accuracy of 97.3\% and 98.6\%, respectively. For property
ownership, we can identify 4.9\% data samples with an accuracy of 99.0\%. For
vehicle ownership, we can identify 14.7\% data samples with an accuracy of
91.9\%. For marital status, we can identify 19.9\% data samples with an accuracy
of 93.5\%. For parental status, we can identify 25.8\% data samples with an
accuracy of 91.3\%.

\revised{Besides, our analysis shows that the distribution of users with high prediction accuracy is demographically representative of the overall collected dataset, indicating no bias toward specific groups. We perform a chi-square test to check consistency. The results show that the distributions of high-confidence data (p-value = 0.87) and other data (p-value = 0.96) align with the distribution of the entire test dataset, further confirming the absence of bias.}

\noindent \textbf{Correlation between Confidence and Accuracy.}
From the red lines in Figure~\ref{fig:results_all} (b)-(h), we can also observe
that \textit{the confidence score produced by our model is positively correlated
with accuracy}. We calculate the Pearson's correlation coefficient between confidence and accuracy, which is $0.992$ on average. It implies that for a given data sample, the higher the
confidence score, the more likely the prediction is to be correct. This is an
important indicator of the effectiveness of the attack. It means that we can
effectively select the subset of users by model confidence, and for the selected
users, we can accurately infer their privacy attributes. 


\subsection{Insights From Data}
\label{sec:exp_insights}

To further understand the reasons for the privacy leakage, we
conduct a thorough investigation into the users that the model has high
confidence scores to identify the sources of leakage for each
attribute. 
We found that the users' privacy attributes relate to the preferred mini-app types and how they operate smartphones.
% 
Specifically, we mainly study three metrics: two metrics are related to \mini: the number of mini-apps (\#Mini-app) and the access frequencies (\#Access) of specific mini-app types, and one metric relates to \op: the number of button clicks (\#Click) in each \op sample.


\noindent \textbf{Gender.}
We display \#Mini-app, \#Access of representative mini-apps in Figure~\ref{fig:gender}~(a) and (b). In
Figure~\ref{fig:gender} (c), we plot the user distribution w.r.t. \#Click. For all figures, the red bars represent female
users, and the blue bars represent male users.
% 
For \mini, we
found that female users prefer \texttt{Shopping} and
\texttt{Beauty} mini-apps. Female users averagely use $6.8$ and
$5.9$ kinds of mini-apps for \texttt{Shopping} and \texttt{Beauty}, with a frequency
of $14.2$ and $10.9$ times. Conversely, male users only access
$2.1$ \texttt{Shopping} mini-apps for $2.4$ times, and $0.5$ \texttt{Beauty} mini-apps
for $0.3$ times. 
% 
We found that male users prefer \texttt{News} ($5.4$ mini-apps for $9.8$ times) and
\texttt{Sports} ($5.8$ mini-apps for $10.6$ times). Conversely, female users
only access $0.7$ mini-apps of these two categories for $0.4$ times.
% 
For \op, female users usually operate more frequently than male
users, confirmed by a recent study~\cite{tian2021and}.
Figure~\ref{fig:gender}~(c) shows that male's \#Click ranges from 10 to
20, while female's \#Click ranges from 20 to 35. This observation aligns
with recent research showing that females tend to use apps for a longer time than
males~\cite{tian2021and}. 







\begin{figure}[!t]
    \centering
    \includegraphics[width=1.\columnwidth]{fig-gender.pdf}
    \vspace{-2em}
	\caption{Detailed analysis on the privacy of \textit{gender}.}
	\label{fig:gender}
 %\vspace{-1em}
\end{figure}



\noindent \textbf{Location.} 
\mini and \op can reflect the economic and population status of
the city where users live, thus leaking the location attribute. To display
this result, we first show the statistical observation of \mini, then map the
mini-app access frequencies to the users' geographical locations in
China, and display the statistics of \#Click.



For \mini, both \#Mini-app and \#Access correlate with users' location. Users in Tier-1, Tier-2, and Tier-3 cities averagely utilize 27.9, 23.6, and 17.3 types of mini-apps. For \#Access, Users in Tier-1, Tier-2, and Tier-3 cities use mini-apps for 36.3, 20.9, and 13.1 times, respectively.
% 
To further verify our observation, we map \#Access to the
users' geographical locations in China (by province) in
Figure~\ref{fig:loc_map}~(a). The darker color represents a higher \#Access. Red and yellow dots represent the Tier-1 and Tier-2 cities. We plot the Heihe-Tengchong Line~\cite{guo2014scientific} with a dashed red line. Heihe-Tengchong Line is a
famous geographical line in China representing population density. The line's southeast side has 94\% of China's population and represents a higher economic level. The northwest side only has 6\% of the population and represents a lower economic level.
From Figure~\ref{fig:loc_map}~(a), we can see that the southeast (lower right in the figure) area of the line is darker than the northwest (upper left) area, which means \#Access is positively correlated with the economic level. We think it is because people in more developed cities rely more on digital services daily. For example, people in Tier-1 cities are more likely to use mini-apps to order food, buy tickets, and pay utility fees.



For \op, we observe a remarkable difference in \#Click between users from different locations. 
Specifically, we focus on the button of \texttt{password-free payment}~\cite{liu2015role}.
This button lets users execute small-amount transactions (below 200 RMB, approximately 30 USD) without inputting passwords. This button is designed to improve user experience and avoid frequently entering passwords. We found that users in the more developed areas have higher \#Click on \texttt{password-free payment}. For users from Tier-3, Tier-2, and Tier-1 cities, \#Click on \texttt{password-free payment} is $0.1$, $0.7$, and $1.1$, respectively. We also observe that this statistical result is specific to the \texttt{password-free payment} button. We do not observe this phenomenon on other buttons with similar functions, such as the \texttt{payment} button. 



\begin{figure}[!t]
    \centering
    \includegraphics[width=0.95\columnwidth]{fig-location.pdf}
    %\vspace{-1em}
	\caption{\revised{Detailed analysis on the privacy of \textit{location}.}}
	\label{fig:loc_map}
 %\vspace{-1em}
\end{figure}






\noindent \textbf{Age.} 
The good performance of \tool for age is because
people of different ages display different patterns on the types of mini-apps and
how they interact with mini-apps. 

For \mini, we plot the number of accessed mini-app categories (\#Category) in
Figure~\ref{fig:age}~(a). Minors (\texttt{Under 18}) and the elderly
(\texttt{Above 65}) access a relatively smaller \#Category (13.2 and 16.9) than other age
groups (above 21.0). 
% 
Notably, minors rarely access mini-apps of \texttt{Finance} because regulations restrict financial services to minors. Similarly, elderly users rarely access mini-apps of 
\texttt{Dating} and \texttt{Comics}, as
older people are not the target users of these mini-apps.






\begin{figure}[!t]
    \centering
    \includegraphics[width=1.\columnwidth]{fig-age.pdf}
    \vspace{-2em}
	\caption{Detailed analysis on the privacy of \textit{age}.}
	\label{fig:age}
%\vspace{-1em}
\end{figure}


For \op, we display \#Click on any buttons (\#Button Click), \#Click on the
\texttt{Payment} button (\#Payment Button Click), and \#Click on
the \texttt{Back} button (\#Back Button Click) of each age group in
Figure~\ref{fig:age}~(b)-(d). 
% 
For \#Button Click, elderly users have the lowest value (average $17.8$) compared to other groups ($31.5$).
This is because the elderly are less familiar with using super-apps and tend to react more slowly to the response.
% 
For \#Payment Button Click, minors rarely click this button (average $0.2$ clicks) than other groups ($1.4$ clicks) because minors are not allowed to make payments without the consent of their guardians. 
% 
For \#Back Button Click, the elderly have the highest value ($3.5$) compared to other groups ($2.3$) because the elderly usually cannot accurately select the desired button and thus need to click the \texttt{Back} button to return to the previous page and try again.






\begin{table}[t]
\centering
\caption{Comparison for users with and without property.}
\vspace{-0.5em}
\label{tab:property}
\begin{adjustbox}{max width = 0.9\columnwidth}
\begin{tabular}{@{}cccc@{}}
\hline
\textbf{Mini-app} & \textbf{Prodived} & \multicolumn{2}{c}{\textbf{\#Access}} \\ \cline{3-4} 
\textbf{Name} & \textbf{Services} & \textbf{W/o Prop.} & \textbf{W/ Prop.} \\ \hline
Utilities & Household bills payment & 0.1 & 2.6 \\
State Grid & Electricity bill payment & 0.0 & 1.3 \\
CSGrid & Electricity bill payment & 0.0 & 1.1 \\
Ziroom & Housing rental & 1.2 & 0.0 \\
Anjuke & Housing rental & 1.1 & 0.0 \\
Renting & Housing rental & 1.6 & 0.0 \\ \hline
\end{tabular}
\end{adjustbox}
%\vspace{-1em}
\end{table}

\noindent \textbf{Property.} 
Table~\ref{tab:property} shows representative mini-apps, the services provided by each mini-app, and the access
frequencies of users with and without property. Users without any
property rarely access mini-apps that can pay household bills (\eg \texttt{Utilities}, \texttt{State Grid},
and \texttt{CSGrid}; frequency below $0.1$). Contrarily, users with property access these mini-apps more frequently (an average of $1.7$ times). This is because
these users must pay the household bills themselves.
% 
On the other side, users without property prefer mini-apps about
house renting (\eg \texttt{Ziroom}, \texttt{Anjuke}, and \texttt{Renting}). The access frequency ($1.3$) is higher than users with property ($0.0$). This is because users without property are likelier to rent a
house.




\begin{table}[t]
\centering
\caption{Comparison for users with and without vehicles.}
\label{tab:vehicle}
\begin{adjustbox}{max width = 0.9\columnwidth}
\begin{tabular}{@{}cccc@{}}
\hline
\textbf{Mini-app} & \textbf{Prodived} & \multicolumn{2}{c}{\textbf{\#Access}} \\ \cline{3-4} 
\textbf{Name} & \textbf{Services} & \textbf{W/o Veh.} & \textbf{W/ Veh.} \\ \hline
12123 & Traffic police platform & 0.2 & 1.5 \\
Sinopec & Fueling & 0.0 & 3.1 \\
CNPC & Fueling & 0.0 & 2.7 \\
Shell & Fueling & 0.0 & 2.3 \\
Car Life & Vehicle insurance, highway toll & 0.0 & 2.8 \\
Halo & Shared bicycle & 12.5 & 2.6 \\
Didi & Taxi & 8.1 & 3.3 \\
Caocao & Taxi & 6.3 & 1.3 \\
Transport & Public transportation & 13.9 & 2.4 \\
Outgoing & Taxi, public transportation & 16.8 & 3.1 \\ \hline
\end{tabular}
\end{adjustbox}
%\vspace{-1em}
\end{table}

\begin{figure}[!t]
    \centering
    \includegraphics[width=1.\columnwidth]{fig-vehicle.pdf}
    \vspace{-2em}
	\caption{Detailed analysis on the privacy of \textit{vehicle}. }
	\label{fig:vehicle}
%\vspace{-1em}
\end{figure}

\noindent \textbf{Vehicle.} 
Table~\ref{tab:vehicle} shows representative mini-apps that may expose vehicle ownership, the services provided by each mini-app, and the access frequencies of users with and without
vehicles. Users without vehicles rarely access mini-apps for traffic police,
fueling, and vehicle insurance (\eg \texttt{12123}, \texttt{Sinopec},
and \texttt{Car Life}; frequency lower than $0.2$). Contrarily, vehicle users access these mini-apps more frequently (average $2.5$ times). 
% 
Users without vehicles prefer public
transportation mini-apps (\eg \texttt{Transport}, \texttt{Didi} and
\texttt{Caocao}). 
In Table~\ref{tab:vehicle}, \#Access for these users ($11.5$) is significantly higher than vehicle users ($2.5$).
This is because users
without vehicles rely on public transportation or a taxi to travel, while
users with vehicles can drive their cars.



For \op, we show \#Click on two buttons in Figure~\ref{fig:vehicle}: \texttt{Back} and \texttt{Password-free
Payment}. The blue and red bars represent the users with and without vehicles. Users
with vehicles have a high \#Clicks on both buttons, higher than users without
vehicle by an average of 38.55\%.
This may be because users
with vehicles are more likely to use super-apps while
driving~\cite{oviedo2020frustrating}, making password-free payments more
convenient and more frequently using the \texttt{Back} button to return to the
previous page. 




\begin{figure}[!t]
    \centering
    \includegraphics[width=1.0\columnwidth]{fig-martial.pdf}
     \vspace{-2em}
	\caption{\revised{Detailed analysis on the privacy of \textit{marital}.}}
	\label{fig:marital}
 %\vspace{-1em}
\end{figure}

\noindent \textbf{Marital.}
We display \#Mini-app and \#Access for representative mini-app categories in
Figure~\ref{fig:marital}. The married users prefer mini-apps about \texttt{Education}, \texttt{House
and home}, \texttt{Finance}, and
\texttt{Medical}. For these types in Figure~\ref{fig:marital}, the red bars (married users) are significantly higher than
the blue bars (unmarried users) by $4.7\times$ for \#Mini-app and $6.3\times$ for \#Access. This means married users focus more on family
aspects of life. 
% 
Conversely, unmarried users prefer mini-apps about
\texttt{Dating} and \texttt{Social}, with $1.5\times$ higher \#Mini-app and $82.0\times$ higher \#Access than married users.





\begin{figure}[!t]
    \centering
    \includegraphics[width=1.0\columnwidth]{fig-parental.pdf}
     \vspace{-2em}
	\caption{Detailed analysis on the privacy of \textit{parental}.}
	\label{fig:parental}
%\vspace{-1em}
\end{figure}

\noindent \textbf{Parental.} 
We display \#Mini-app and \#Access for representative
mini-app categories in Figure~\ref{fig:parental}~(a) and (b), and show \#ButtonClick in Figure~\ref{fig:parental}~(c).
% 
For \mini, users with children prefer the mini-apps about
\texttt{Education}, \texttt{Parenting}, and
\texttt{Books}. In Figure~\ref{fig:parental}~(a) and
(b), the \#Mini-app and \#Access of users with children (red bars) are
higher than those without children (blue bars) by $9.6\times$ and $10.0\times$,
respectively. We believe the reason is that parents tend to pay more attention
to their children's education.
% 
For \op, users with children have a lower \#ButtonClick than users
without children. In Figure~\ref{fig:parental}~(c), the \#ButtonClick of
parents (red bar) is $18.0\%$ lower than childless users (blue bar). 



\subsection{Architecture and Calibration Comparison}

\label{sec:architecture_comparison}

In this section, we study different choices of model architecture and confidence
calibration techniques. For the model architecture, we select three
representative architectures: a CNN-based model, an RNN-based model, and our
Transformer-based model. The CNN-based model is a ResNet model with 48
convolutional layers~\cite{lecun1995convolutional, he2016deep}. The RNN-based
model has 8 LSTM layers following the default architecture of prior
literature~\cite{salehinejad2017recent, hochreiter1997long}. All
architectures have been demonstrated to be effective to extract high-dimensional
information from time series
data~\cite{zerveas2021transformer,hong2020holmes,ma2021effectiveness,ding2019modeling}.
For calibration, we compare the temperature calibration with two
representative techniques: vector calibration and matrix
calibration~\cite{pmlr-v70-guo17a}. 



\noindent \textbf{Evaluation Metrics.}
We use four different metrics to evaluate the performance: the proportion of high-confidence data (PHC), precision,
recall, and F1-score. 
% 
Given a confidence threshold, the PHC measures the proportion of data samples that
the model can confidently infer. A higher PHC indicates a higher attack performance.
% 
Precision, recall, and F1-score are three widely used metrics to evaluate classification models. We use these three metrics to
comprehensively evaluate the model's correctness on the PHC data samples.


\noindent \textbf{Results.}
We display the averaged results in Table~\ref{tab:results}, in which the values are
averaged across the seven labels of privacy attributes.
% 
First, for all settings, the PHC is above $8.8\%$. Precision, recall, and F1-score
are all above $89.6\%$. The
proposed attack is general to different architectures
and calibration techniques. 
% 
Second, among all the settings, the Transformer-based model with temperature
calibration achieves the best performance. The PHC is $16.1\%$, and the
precision, recall, and F1-score are all above $95.4\%$. This means that this setting can
correctly infer the privacy of the most number of users with the highest
accuracy. 



\newcolumntype{Y}{>{\centering\arraybackslash}X}
\begin{table*}[]
\centering
\caption{Comparison between model architectures and calibration techniques. For each metric, we report the average value across seven privacy attributes.}
%\vspace{-1em}
\label{tab:results}
\begin{adjustbox}{max width = 0.9\linewidth}
\begin{tabularx}{\linewidth}{Y|YYY|YYY|YYY}
\hline
 Model & \multicolumn{3}{c|}{CNN-based}  & \multicolumn{3}{c|}{RNN-based} & \multicolumn{3}{c}{Transformer-based} \\ \hline
 Calibration & Temperature & Vector & Matrix & Temperature & Vector & Matrix & Temperature & Vector & Matrix \\ \hline
 PHC & 14.6\%  & 11.8\% & 15.9\% & 10.0\% & 7.9\% & 9.8\%  & \textbf{16.1\%} & 15.6\% & 8.8\% \\
 Precision & 91.9\% & 91.3\% & 90.4\% & 91.1\% & 92.2\% & 91.4\% & \textbf{95.5\%} & 94.4\% & 95.3\% \\
 Recall & 91.8\% & 90.6\% & 89.8\% & 90.2\% & 90.5\% & 89.6\% & \textbf{95.4\%} & 94.1\% & 94.4\% \\
 F1-score & 91.8\% & 90.8\% & 89.9\% & 90.5\% & 91.1\% & 90.1\% & \textbf{95.4\%} & 94.2\% & 94.8\% \\ \hline
\end{tabularx}
\end{adjustbox}
\end{table*}

\subsection{Ablation Study}
\label{sec:ablation_study}
In this section, we compare the results of different settings used in \tool.
\begin{figure}[t]
    \centering
    \includegraphics[width=1.0\columnwidth]{fig-ablation1.pdf}
    \vspace{-2em}
	\caption{Comparision of different settings.}
	\label{fig:appendixsettings}
 \vspace{-1em}
\end{figure}

\noindent \textbf{Threshold.}
First, we assessed the impact of the threshold settings. A higher threshold means we only attack samples with higher confidence scores, implying a smaller victim subset and higher accuracy. Specifically, we compared the proportion of identified samples and their accuracy when setting the threshold at 0.9, 0.8, and 0.7. In Figure~\ref{fig:appendixsettings} (a), light-colored bars represent the proportion of samples that can be identified, and dark-colored bars represent the proportion of accurately inferred samples. We can observe that a lower threshold encompasses more samples, but the accuracy is relatively decreased. The results for the three threshold settings are displayed in Figure~\ref{fig:appendixsettings} (b), where a 0.9 threshold achieves an average accuracy of 95.7\%, while 0.8 and 0.7 only reach average accuracies of 88.7\% and 82.8\% respectively. \revised{These findings demonstrate that higher confidence scores align with higher inference accuracy, thereby validating our model calibration technique. }In this paper, we therefore adopt a threshold of 0.9 to maximize our attack performance.


\noindent \textbf{Input Data.}
\tool's input data is a concatenation of \mini and \op. To further demonstrate that both data types are privacy-related, we trained models using \mini and \op separately and reported the results in Figure~\ref{fig:appendixsettings} (c). Only using \mini or \op achieves an accuracy of 61.1\% and 54.4\%, respectively. Although the values are lower than \tool's concatenation solution (65.2\%), they are higher than the random guess baseline by an average of 44.0\%. Thus, both \mini and \op are effective for \tool.


\noindent \textbf{Data Length.}
In \tool, we empirically set $N=200$. In this section, we compare it with $N=100$ and $N=400$. As shown in Figure~\ref{fig:appendixsettings} (d), using 
$N=400$ only provides a marginal improvement of 0.2\%, while $N=100$ decreases accuracy by 5.4\%. Therefore, our setting is an effective choice.

